\documentclass[10pt,a4paper]{amsart}
\usepackage{amsmath,amssymb,amsthm,graphicx}
% Retained preamble for the Stacks Project chapter snapshot in this directory.
% The chapter's own preamble.tex is a copy of the shared Stacks Project preamble; the block
% below is sliced MECHANICALLY, line for line, from it. Nothing was added, removed or
% redefined; the wrapper class, the xr-hyper/hyperref wiring and the externaldocument list
% are omitted because the assembler supplies the document class.
%
% lines taken from preamble.tex: 3, 5, 9, 11, 12, 156, 157, 158, ... (51 lines total)
%
% package coverage restored from the chapter's own preamble (2026-10-13): xy, xy-2cell

\usepackage[all]{xy}
\xyoption{2cell}
\UseAllTwocells

\newenvironment{reference}{\comment}{\endcomment}
\newenvironment{slogan}{\comment}{\endcomment}
\newenvironment{history}{\comment}{\endcomment}
\theoremstyle{plain}
\newtheorem{theorem}[subsection]{Theorem}
\newtheorem{proposition}[subsection]{Proposition}
\newtheorem{lemma}[subsection]{Lemma}
\theoremstyle{definition}
\newtheorem{definition}[subsection]{Definition}
\newtheorem{example}[subsection]{Example}
\newtheorem{exercise}[subsection]{Exercise}
\newtheorem{situation}[subsection]{Situation}
\theoremstyle{remark}
\newtheorem{remark}[subsection]{Remark}
\newtheorem{remarks}[subsection]{Remarks}
\numberwithin{equation}{subsection}
\def\lim{\mathop{\mathrm{lim}}\nolimits}
\def\colim{\mathop{\mathrm{colim}}\nolimits}
\def\Spec{\mathop{\mathrm{Spec}}}
\def\Hom{\mathop{\mathrm{Hom}}\nolimits}
\def\Ext{\mathop{\mathrm{Ext}}\nolimits}
\def\SheafHom{\mathop{\mathcal{H}\!\mathit{om}}\nolimits}
\def\SheafExt{\mathop{\mathcal{E}\!\mathit{xt}}\nolimits}
\def\Sch{\mathit{Sch}}
\def\Mor{\mathop{\mathrm{Mor}}\nolimits}
\def\Ob{\mathop{\mathrm{Ob}}\nolimits}
\def\Sh{\mathop{\mathit{Sh}}\nolimits}
\def\NL{\mathop{N\!L}\nolimits}
\def\CH{\mathop{\mathrm{CH}}\nolimits}
\def\proetale{{pro\text{-}\acute{e}tale}}
\def\etale{{\acute{e}tale}}
\def\QCoh{\mathit{QCoh}}
\def\Ker{\mathop{\mathrm{Ker}}}
\def\Im{\mathop{\mathrm{Im}}}
\def\Coker{\mathop{\mathrm{Coker}}}
\def\Coim{\mathop{\mathrm{Coim}}}
\DeclareMathSymbol{\boxtimes}{\mathbin}{AMSa}{"02}
\def\QCohstack{\mathcal{QC}\!\mathit{oh}}
\def\Cohstack{\mathcal{C}\!\mathit{oh}}
\def\Spacesstack{\mathcal{S}\!\mathit{paces}}
\def\Quotfunctor{\mathrm{Quot}}
\def\Hilbfunctor{\mathrm{Hilb}}
\def\Curvesstack{\mathcal{C}\!\mathit{urves}}
\def\Polarizedstack{\mathcal{P}\!\mathit{olarized}}
\def\Complexesstack{\mathcal{C}\!\mathit{omplexes}}
\def\Pic{\mathop{\mathrm{Pic}}\nolimits}
\def\Picardstack{\mathcal{P}\!\mathit{ic}}
\def\Picardfunctor{\mathrm{Pic}}
\def\Deformationcategory{\mathcal{D}\!\mathit{ef}}
\usepackage{cleveref}
\crefname{originallemma}{original lemma}{original lemmas}
\Crefname{originallemma}{Original lemma}{Original lemmas}
\makeatletter
\newlabel{coherent-lemma-algebraize-morphism}{{0A42}{source}{}{sourcepointer-1}{}}
\newlabel{coherent-lemma-algebraize-morphism@cref}{{[originallemma][0][]0A42}{[1][0][]source}{}{}{}}
\makeatother
\crefname{originallemma}{original lemma}{original lemmas}
\Crefname{originallemma}{Original lemma}{Original lemmas}
\makeatletter
\newlabel{examples-defos-lemma-schemes-morphisms-smooth-to-base}{{0E3X}{source}{}{sourcepointer-2}{}}
\newlabel{examples-defos-lemma-schemes-morphisms-smooth-to-base@cref}{{[originallemma][0][]0E3X}{[1][0][]source}{}{}{}}
\makeatother
\crefname{originallemma}{original lemma}{original lemmas}
\Crefname{originallemma}{Original lemma}{Original lemmas}
\makeatletter
\newlabel{lemma-curves-qs-lfp}{{0DSS}{source}{}{sourcepointer-3}{}}
\newlabel{lemma-curves-qs-lfp@cref}{{[originallemma][0][]0DSS}{[1][0][]source}{}{}{}}
\makeatother
\crefname{originallemma}{original lemma}{original lemmas}
\Crefname{originallemma}{Original lemma}{Original lemmas}
\makeatletter
\newlabel{more-algebra-section-permanence-completion}{{07NU}{source}{}{sourcepointer-4}{}}
\newlabel{more-algebra-section-permanence-completion@cref}{{[originallemma][0][]07NU}{[1][0][]source}{}{}{}}
\makeatother
\crefname{originallemma}{original lemma}{original lemmas}
\Crefname{originallemma}{Original lemma}{Original lemmas}
\makeatletter
\newlabel{spaces-limits-lemma-descend-finite-presentation}{{07SK}{source}{}{sourcepointer-5}{}}
\newlabel{spaces-limits-lemma-descend-finite-presentation@cref}{{[originallemma][0][]07SK}{[1][0][]source}{}{}{}}
\makeatother
\crefname{originallemma}{original lemma}{original lemmas}
\Crefname{originallemma}{Original lemma}{Original lemmas}
\makeatletter
\newlabel{spaces-more-morphisms-lemma-projective-over-complete}{{0AE7}{source}{}{sourcepointer-6}{}}
\newlabel{spaces-more-morphisms-lemma-projective-over-complete@cref}{{[originallemma][0][]0AE7}{[1][0][]source}{}{}{}}
\makeatother
\crefname{originallemma}{original lemma}{original lemmas}
\Crefname{originallemma}{Original lemma}{Original lemmas}
\makeatletter
\newlabel{spaces-morphisms-lemma-flat-base-change-scheme-theoretic-image}{{089E}{source}{}{sourcepointer-7}{}}
\newlabel{spaces-morphisms-lemma-flat-base-change-scheme-theoretic-image@cref}{{[originallemma][0][]089E}{[1][0][]source}{}{}{}}
\makeatother
\crefname{originallemma}{original lemma}{original lemmas}
\Crefname{originallemma}{Original lemma}{Original lemmas}
\makeatletter
\newlabel{stacks-limits-lemma-representable-by-spaces-limit-preserving}{{0CMV}{source}{}{sourcepointer-8}{}}
\newlabel{stacks-limits-lemma-representable-by-spaces-limit-preserving@cref}{{[originallemma][0][]0CMV}{[1][0][]source}{}{}{}}
\makeatother
\title{Stacks, Moduli of Curves: lemma-contract-basic-uniqueness}
\author{The Stacks Project Authors}
\date{stacks-project:moduli-curves.tex:local-snapshot-20261012; GNU FDL 1.2 or later}
\begin{document}
\maketitle
\noindent\textit{Editorial note.} Statement span is lines 2899-2912 of sources/stacks-moduli-curves/pinned-original-source.tex; proof span is lines 2914-2984, the complete author proof (69 body lines).\par\medskip
\noindent\textit{Editorial note.} Navigation stubs: 8 cited-but-not-included labels are supplied as non-load-bearing pointers with their pinned tags where the snapshot tag index has them.\par\medskip
\noindent\textit{Editorial note.} extra\_packages is empty on purpose: the retained preamble already loads xy with options and a second declaration clashes in the assembler wrapper.\par\medskip
\noindent\textit{Editorial note.} Text is retained verbatim, including the author's own typos and placeholders.\par\medskip
\noindent\textit{Editorial note.} Queue index 4 of queue\_moduli-curves.json (pb=3316, stmt 2899-2912, proof 2914-2984).\par\medskip
\noindent\textit{Allowed standard prerequisites (editorial).}\begin{itemize}
\item Explicit prerequisite boundary. The cited inputs are prerequisites; the argument in the retained span is not replaced by any cited result.
\end{itemize}\medskip
\ifdefined\hypertarget\hypertarget{sourcepointer-1}{}\fi
\noindent\textit{Source reference.} Original lemma 0A42: sources/stacks-moduli-curves/pinned-original-source.tex, source lines 2984--2984. Cited by the retained statement or proof but not part of the collected spans; supplied only as a navigation pointer. The argument does not depend on its content.\par\medskip
\ifdefined\hypertarget\hypertarget{sourcepointer-2}{}\fi
\noindent\textit{Source reference.} Original lemma 0E3X: sources/stacks-moduli-curves/pinned-original-source.tex, source lines 2984--2984. Cited by the retained statement or proof but not part of the collected spans; supplied only as a navigation pointer. The argument does not depend on its content.\par\medskip
\ifdefined\hypertarget\hypertarget{sourcepointer-3}{}\fi
\noindent\textit{Source reference.} Original lemma 0DSS: sources/stacks-moduli-curves/pinned-original-source.tex, source lines 655--658. Cited by the retained statement or proof but not part of the collected spans; supplied only as a navigation pointer. The argument does not depend on its content.\par\medskip
\ifdefined\hypertarget\hypertarget{sourcepointer-4}{}\fi
\noindent\textit{Source reference.} Original lemma 07NU: sources/stacks-moduli-curves/pinned-original-source.tex, source lines 2984--2984. Cited by the retained statement or proof but not part of the collected spans; supplied only as a navigation pointer. The argument does not depend on its content.\par\medskip
\ifdefined\hypertarget\hypertarget{sourcepointer-5}{}\fi
\noindent\textit{Source reference.} Original lemma 07SK: sources/stacks-moduli-curves/pinned-original-source.tex, source lines 2984--2984. Cited by the retained statement or proof but not part of the collected spans; supplied only as a navigation pointer. The argument does not depend on its content.\par\medskip
\ifdefined\hypertarget\hypertarget{sourcepointer-6}{}\fi
\noindent\textit{Source reference.} Original lemma 0AE7: sources/stacks-moduli-curves/pinned-original-source.tex, source lines 2984--2984. Cited by the retained statement or proof but not part of the collected spans; supplied only as a navigation pointer. The argument does not depend on its content.\par\medskip
\ifdefined\hypertarget\hypertarget{sourcepointer-7}{}\fi
\noindent\textit{Source reference.} Original lemma 089E: sources/stacks-moduli-curves/pinned-original-source.tex, source lines 2984--2984. Cited by the retained statement or proof but not part of the collected spans; supplied only as a navigation pointer. The argument does not depend on its content.\par\medskip
\ifdefined\hypertarget\hypertarget{sourcepointer-8}{}\fi
\noindent\textit{Source reference.} Original lemma 0CMV: sources/stacks-moduli-curves/pinned-original-source.tex, source lines 2984--2984. Cited by the retained statement or proof but not part of the collected spans; supplied only as a navigation pointer. The argument does not depend on its content.\par\medskip
\begin{lemma}
\label{lemma-contract-basic-uniqueness}
Let $S$ be a scheme and $s \in S$ a point.
Let $f : X \to S$ and $g_i : Y_i \to S$, $i = 1, 2$ be families of curves.
Let $c_i : X \to Y_i$ be morphisms over $S$.
Assume there is an isomorphism $Y_{1, s} \cong Y_{2, s}$
of fibres compatible with $c_{1, s}$ and $c_{2, s}$.
If $c_{1, s, *}\mathcal{O}_{X_s} = \mathcal{O}_{Y_{1, s}}$ and
$R^1c_{1, s, *}\mathcal{O}_{X_s} = 0$, then there exist an
open neighbourhood $U$ of $s$ and an isomorphism
$Y_{1, U} \cong Y_{2, U}$ of families of curves over $U$
compatible with the given isomorphism of fibres and with
$c_1$ and $c_2$.
\end{lemma}

\begin{proof}
Recall that $\mathcal{O}_{S, s} = \colim \mathcal{O}_S(U)$ where
the colimit is over the system of affine neighbourhoods $U$ of $s$.
Thus the category of algebraic spaces of finite presentation over
the local ring is the colimit of the categories of algebraic spaces
of finite presentation over the affine neighbourhoods of $s$.
See Limits of Spaces, Lemma
\ref{spaces-limits-lemma-descend-finite-presentation}.
In this way we reduce to the case where $S$ is the spectrum of
a local ring and $s$ is the closed point.

\medskip\noindent
Assume $S = \Spec(A)$ where $A$ is a local ring and $s$ is the closed point.
Write $A = \colim A_j$ with $A_j$ local Noetherian (say essentially
of finite type over $\mathbf{Z}$) and local transition homomorphisms.
Set $S_j = \Spec(A_j)$ with closed point $s_j$. We can find a
$j$ and families of curves $X_j \to S_j$, $Y_{j, i} \to S_j$,
see Lemma \ref{lemma-curves-qs-lfp} and
Limits of Stacks, Lemma
\ref{stacks-limits-lemma-representable-by-spaces-limit-preserving}.
After possibly increasing $j$ we can find morphisms
$c_{j, i} : X_j \to Y_{j, i}$ whose base change to $s$ is $c_i$, see
Limits of Spaces, Lemma
\ref{spaces-limits-lemma-descend-finite-presentation}.
Since $\kappa(s) = \colim \kappa(s_j)$ we can similarly
assume there is an isomorphism $Y_{j, 1, s_j} \cong Y_{j, 2, s_j}$
compatible with $c_{j, 1, s_j}$ and $c_{j, 2, s_j}$.
Finally, the assumptions
$c_{1, s, *}\mathcal{O}_{X_s} = \mathcal{O}_{Y_{1, s}}$ and
$R^1c_{1, s, *}\mathcal{O}_{X_s} = 0$
are inherited by $c_{j, 1, s_j}$ because
$\{s_j \to s\}$ is an fpqc covering and
$c_{1, s}$ is the base of $c_{j, 1, s_j}$ by this covering (details omitted).
In this way we reduce the lemma to the case discussed in the next paragraph.

\medskip\noindent
Assume $S$ is the spectrum of a Noetherian local ring $\Lambda$ and
$s$ is the closed point. Consider the scheme theoretic image $Z$ of
$$
(c_1, c_2) : X \longrightarrow Y_1 \times_S Y_2
$$
The statement of the lemma is equivalent to the assertion that $Z$ maps
isomorphically to $Y_1$ and $Y_2$ via the projection morphisms.
Since taking the scheme theoretic image of this morphism commutes
with flat base change (Morphisms of Spaces, Lemma
\ref{spaces-morphisms-lemma-flat-base-change-scheme-theoretic-image},
we may replace $\Lambda$ by its completion
(More on Algebra, Section \ref{more-algebra-section-permanence-completion}).

\medskip\noindent
Assume $S$ is the spectrum of a complete Noetherian local ring $\Lambda$.
Observe that $X$, $Y_1$, $Y_2$ are schemes in this case
(More on Morphisms of Spaces, Lemma
\ref{spaces-more-morphisms-lemma-projective-over-complete}).
Denote $X_n$, $Y_{1, n}$, $Y_{2, n}$ the base changes of
$X$, $Y_1$, $Y_2$ to $\Spec(\Lambda/\mathfrak m^{n + 1})$.
Recall that the arrow
$$
\Deformationcategory_{X_s \to Y_{2, s}} \cong
\Deformationcategory_{X_s \to Y_{1, s}} \longrightarrow
\Deformationcategory_{X_s}
$$
is an equivalence, see Deformation Problems, Lemma
\ref{examples-defos-lemma-schemes-morphisms-smooth-to-base}.
Thus there is an isomorphism of formal objects
$(X_n \to Y_{1, n}) \cong (X_n \to Y_{2, n})$
of $\Deformationcategory_{X_s \to Y_{1, s}}$.
Finally, by Grothendieck's algebraization theorem
(Cohomology of Schemes, Lemma \ref{coherent-lemma-algebraize-morphism})
this produces an isomorphism $Y_1 \to Y_2$ compatible with $c_1$ and $c_2$.
\end{proof}

\end{document}
